 
%% ============================================================
%%  General Introduction
%% ============================================================
\chapter*{General Introduction}
\addcontentsline{toc}{chapter}{General Introduction}

Since the dawn of civilization, knowledge transmission has been one of the fundamental
pillars of human progress. In today's digital era, educational systems are facing new
challenges, including the globalization of skills, the diversification of learner profiles,
the rapid evolution of technology, and the growing demand for accessible and flexible
learning opportunities. In this context, e-learning has emerged as an effective solution,
providing learners with personalized, flexible, and accessible educational experiences.

In Tunisia and across the MENA region, the education and training sector is undergoing a
significant digital transformation. Universities, training centers, educational institutions,
and independent instructors are increasingly seeking digital solutions capable of managing
their educational activities while preserving their organizational structures and teaching
methodologies. This growing demand highlights the need for a multi-organizational platform
that enables multiple institutions to operate within a shared infrastructure while maintaining
their autonomy, security, and customization requirements.

It is within this context that the present final-year project was carried out at
\textbf{YohaTech Monastir} under the title:

\begin{center}
\textit{``Development and Implementation of an AI-Powered Multi-Organization Learning Management System''}
\end{center}
Tesla Academy is a Learning Management System (LMS) designed around a multi-organization architecture. The platform enables multiple organizations, such as universities, training centers, educational institutions, and independent instructors, to manage their educational activities through isolated and secure workspaces while sharing a common technological infrastructure.

In addition to traditional LMS functionalities, Tesla Academy incorporates Artificial
Intelligence technologies aimed at enhancing both teaching and learning experiences. The
platform provides intelligent educational assistance, automated question generation,
assignment evaluation and feedback, content retrieval through Retrieval-Augmented Generation
(RAG), and AI-powered support tools that facilitate knowledge acquisition and improve learner
engagement.

The project was developed using a modern software architecture composed of a web-based
frontend application, a backend service responsible for business logic and data management,
and a dedicated AI microservice that delivers advanced intelligent functionalities. This
architecture ensures modularity, maintainability, scalability, and seamless integration of
AI capabilities within the learning environment.

The primary objective of this project is to design, develop, and deploy a secure, scalable,
and intelligent learning platform capable of serving multiple organizations while leveraging
Artificial Intelligence to improve educational effectiveness, learner support, and overall
user experience.

Finally, this report presents the different phases of the project, from the study of the
existing solutions and requirements analysis to system design, implementation, testing, and
deployment. It concludes with a summary of the achievements and proposes future enhancements
for the continuous evolution of Tesla Academy.
 
